\subsubsection{ONT RNA Sample Preparation}

One mL aliquots of actively growing JCVI-syn3A cells were placed in 25 mL of SP4 + 17\% KnockOut serum replacement in 50 mL conical tubes and incubated overnight at 37\textdegree C.
The cells were harvested in mid-late log phase as indicated by a red-orange color of the phenol red pH indicator in the growth media.
The cultures were harvested when the media pH was 7.0.
JCVI-syn3A cultures achieved this growth stage in 22 hours.
Growth was arrested in the bulk cultures by placing the tubes on wet ice for 10 minutes.
Keeping all tubes on ice, 1.4 mL of bulk cultures were distributed into fifteen 1.8 mL microfuge tubes per culture.
These were centrifuged at 4\textdegree C for five minutes at 10,000$\times$g using a tabletop refrigerated centrifuge to pellet the cells.
The supernatants from ten of the fifteen tubes were aspirated.
In the remaining five tubes, the supernatant was then used to suspend pellets.
The cell suspensions from those five tubes were then used to resuspend the pellets in five other tubes, and those more concentrated cell suspensions were used to resuspend the cells in the last five tubes.
The JCVI-syn3A cells were now in five tubes containing concentrated cell suspensions.
These tubes were centrifuged for five minutes at 4\textdegree C under 10,000$\times$g with a tabletop refrigerated centrifuge.
All supernatants were aspirated completely, and the remaining pellets were snap frozen using liquid nitrogen.
Sample tubes were transferred to dry ice and immediately stored at -80\textdegree C.

In the next step, cellular RNA was extracted from JCVI-syn3A cell samples using a phenol-chloroform extraction protocol based off procedures adapted from~\cite{grunberger_next_2019} and the ThermoFisher TRIzol\textsuperscript{TM} Reagent user guide~\cite{thermofisher_trizol}.
Briefly, 1 mL of 4\textdegree C TRIzol\textsuperscript{TM} was added to each tube containing frozen cell pellets.
The liquid was pipetted up and down multiple times to break up the cell pellet.
The tubes were incubated at room temperature for 5 minutes to allow complete disassociation of nucleoprotein complexes.
200 $\mu$L of chloroform was then added and the tubes were mixed by shaking for 30 seconds.
After the tubes incubated at room temperature for 3 minutes, they were centrifuged at 12,000$\times$g for 15 minutes at 4\textdegree C.
The upper RNA-containing aqueous phase was pipetted without collecting any of the white interphase separating the aqueous from the phenol-chloroform lower phase, and placed in fresh 1.8 mL tubes already containing 500 $\mu$L isopropanol.
The tubes were incubated on ice for 10 minutes and then centrifuged at 12,000$\times$g for 10 minutes at 4\textdegree C.
The supernatant was aspirated and discarded off the white gel-like pellets.
The tubes were inverted and the pellets were allowed to air dry for 15 minutes.
The RNA pellets were then dissolved in 20 $\mu$L of RNase-free water containing 0.1 mM EDTA.
RNA concentration was assessed with a Qubit RNA HS fluorometer assay and an Agilent TapeStation 2200.
The next step was ribosomal RNA (rRNA) depletion.

The JCVI-syn3A samples were rRNA depleted using the NEBNext\textregistered{} rRNA Depletion Kit for bacteria.
One $\mu$g of RNA was brought up to 11 $\mu$L with RNase-free water.
To this were added 2 $\mu$L of NEBNext\textregistered{} Bacterial rRNA Depletion Solution and 2 $\mu$L of NEBNext\textregistered{} Probe Hybridization Buffer.
The reaction was heated to 95\textdegree C for 2 minutes and cooled to 22\textdegree C at a ramp rate of 0.1\textdegree C/sec.
This was followed immediately by RNase H digestion, where 2 $\mu$L of RNase H Reaction Buffer and 2 $\mu$L of NEBNext Thermostable RNase H were added to the hybridized RNA solution to a total volume of 20 $\mu$L.
This reaction was incubated at 50\textdegree C for 30 minutes.

Following a bead cleanup with Agencourt RNAClean XP beads (Beckman Coulter\texttrademark, cat \# A63987), the 15 $\mu$L of the rRNA solution was Poly(A) tailed at the 3$'$ end.
RNA samples were mixed with a premade solution containing 2 $\mu$L of 10X \eco~Poly(A) Polymerase Reaction Buffer, 2 $\mu$L of ATP and 1 $\mu$L of the \eco~Poly(A) Polymerase (New England Biolabs) and incubated at 37\textdegree C for 30 minutes.
The reaction was stopped by adding EDTA to a final concentration of 10 mM.
Following another bead cleanup along with necessary quality control analysis on the samples, RNA was normalized to 9 $\mu$L containing 50 ng for JCVI-syn3A RNA samples.

\subsubsection{RNA Sample Quality Control Analysis}

For the JCVI-syn3A samples, both a Qubit RNA HS fluorometer assay and an Agilent TapeStation 2200 were used for assessing RNA concentration post-extraction from minimal cells.
Initial Qubit values post RNA extraction were 1890 ng/$\mu$L for JCVI-syn3A.
The Agilent TapeStation 2200 was used throughout the extraction and rRNA depletion steps to assess the quality of the RNA and depletion.
TapeStation results gave an estimate of depletion through the absence of 16S and 23S peaks (Figure SXX).
The corresponding graphs can be found in rRNA Supplemental Data.
A Qubit RNA HS fluorometer was used to detect final quantitation of RNA.

Post poly-adenylation of RNA transcripts, the extent of the DNA and buffer contamination in each of the RNA samples was tested by performing standard spectroscopic measurements (NanoDrop One) and using the Qubit 1$\times$ dsDNA HS assay kit (ThermoFisher Scientific).
Input RNA samples for Oxford Nanopore Direct RNA Sequencing library preparation were finally quantified using the Qubit RNA HS assay kit.
RNA concentration values were used to normalize samples to the desired input volume of 9 $\mu$L containing 50--500 ng of input RNA for library preparation.

\subsubsection{ONT Library Preparation and Direct Sequencing}

After normalization, libraries for Oxford Nanopore sequencing were prepared using the ONT Direct RNA Sequencing Kit (SQK-RNA002) with slight modifications.
Specifically, 0.5 $\mu$L of the RT adapter enzyme (reduced from 1 $\mu$L) and 4 $\mu$L of the RMX RNA adapter (reduced from 6 $\mu$L) were used during library preparation for JCVI-syn3A.
The samples were sequenced using R9.4 flow cells on the ONT GridION platform.
Live base-calling was enabled through the recommended MinKNOW (v23.07.5) scripts to generate fast5 files.
734.08~k reads were generated for JCVI-syn3A over a single run, with an average N50 of 430.

\subsubsection{Processing of Raw ONT Reads}

Results of the Nanopore reads were collected as individual fast5 files.
The individual Nanopore reads were then combined into multi-read files for easier data analysis.
Current signals were converted to canonical bases by base-calling.
For the JCVI-syn3A experiments, base-calling was performed with Guppy (v7.0.9) using the RNA trim strategy and with calibration detection and filter enabled.
In all run reports, base-called bases were 99.99--100\% confident.

\subsubsection{Mapping ONT Reads to Syn3A Genome}


Reads were aligned to the JCVI-syn3A reference genome (NCBI accession CP016816.2) using \texttt{minimap2} v2.30~\cite{li_minimap2_2018} with the long-read Oxford Nanopore preset for non-spliced alignment: \texttt{-ax map-ont -p 0.99 --MD}.
The \texttt{map-ont} preset was selected in preference to the splice-aware preset commonly used for ONT direct RNA because JCVI-syn3A, like its parent \Mmy, is intron-less; allowing N-skip (spliced) CIGAR operations would yield single-read alignments that bridge tens to hundreds of kilobases of unrelated reference.
Strandedness of direct-RNA reads is preserved through the BAM flag bits.
Resulting SAM output was converted, sorted, and indexed using \texttt{samtools} v1.22.1~\cite{li2009samtools}.

Of 734~k input reads, 559~k (76.2\%) mapped as primary alignments to the syn3A genome, with a mean read length of 383~nt (range 49--2{,}858~nt) and mean per-read Phred quality of 31.2; 98.2\% of primary mapped reads exceeded a mean quality of Q20.
The 175~k unmapped reads (23.9\%) were of substantially lower quality (mean Q20.7), with 42.0\% falling below Q15 and 61.0\% shorter than 300~nt, consistent with the expected tail of low-quality and truncated molecules characteristic of ONT direct RNA sequencing rather than reference mismatch.
Secondary alignments accounted for 3.6\% of primary mapped reads (20.2~k alignments) and clustered tightly at two genomic loci ($\sim$55{,}460 and $\sim$343{,}267~bp), corresponding to the two ribosomal RNA operons of syn3A.
A further 7.3~k supplementary alignments (split-read mappings) were retained.

% Reads were aligned to the JCVI-syn3A reference genome (NCBI accession CP016816.2) using \texttt{minimap2} v2.30~\cite{li_minimap2_2018} with the splice-aware preset: \texttt{-ax splice -uf -k14 -p 0.99 --MD}. The \texttt{-uf} flag enforces forward-strand alignment, exploiting the inherent strandedness of direct RNA reads. Resulting SAM output was converted, sorted, and indexed using \texttt{samtools} v1.x~\cite{li2009samtools}.

% Of 734~k input reads, 541~k (73.7\%) mapped as primary alignments
% to the syn3A genome, with a mean read length of 390~nt and mean per-read
% Phred quality of 31.5; 98.4\% of primary mapped reads exceeded a mean
% quality of Q20. The 193~k unmapped reads (26.3\%) were of substantially
% lower quality (mean Q21.0), with 38.2\% falling below Q15 and 65.3\% shorter
% than 300~nt, consistent with the expected tail of low-quality and truncated
% molecules characteristic of ONT direct RNA sequencing rather than reference
% mismatch. Secondary alignments accounted for only 2.7\% of primary mapped
% reads (14.6~k alignments) and clustered tightly
% at two genomic loci ($\sim$55{,}460 and $\sim$343{,}267~bp), corresponding
% to the two ribosomal RNA operons of syn3A.

\subsubsection{Mapping Illumina Reads to Syn3A Genome}

Raw Illumina reads (paired-end, 51 nt, $\sim$12.7~M pairs per replicate) of Syn3A wild type at log-phase from~\cite{sandberg_adaptive_2023} were retrieved from NCBI SRA database and assessed with \texttt{FastQC} and aggregated with \texttt{MultiQC}~\cite{ewels2016multiqc}.
Per-base and per-read quality scores passed all checks and no Illumina adapters were detected, so no trimming was applied.

% The FastQC \emph{sequence duplication} and \emph{per-base sequence content} flags failed in both mates as expected for RNA-seq on a 543~kb bacterial genome (the same transcript fragment is sequenced many times, and the first $\sim$10~nt show the canonical random-hexamer priming bias); neither was treated as a quality issue. A small fraction of overrepresented sequences mapped to the 16S/23S/5S ribosomal operons, indicating residual rRNA after Ribo-Zero depletion, but these were retained for downstream strand-split depth analysis.

Reads were aligned to the JCVI-syn3A reference (NCBI accession CP016816.2) with \texttt{bowtie2} v2.5.5~\cite{langmead2012bowtie2} in default paired-end mode.
The overall alignment rate was 98.88\%: 83.55\% of pairs mapped concordantly exactly once, 10.43\% concordantly with $>$1 hit (corresponding to the two ribosomal RNA operons of syn3A), and 6.02\% of pairs were discordant or unpaired.
SAM output was converted, sorted, and indexed with \texttt{samtools} v1.22.1~\cite{li2009samtools}.
Strand-specific BAMs were generated using the dUTP / fr-firststrand flag convention (read 2 forward on $+$ strand and read 1 reverse on $+$ strand define plus-strand reads; their reverse complements define minus-strand reads), matching the Kapa Stranded RNA-Seq Kit library preparation, and per-strand sequencing-depth bedGraph tracks were computed with \texttt{samtools depth}.

\subsubsection{Quantification of TPM from Sequencing Depth}

Per-gene sense and antisense TPM were computed for~\syna~from the strand-specific depth tracks, using the same depth-based definition applied to the~\syn~Illumina data, namely a gene-length-normalized mean per-base depth scaled to transcripts per million against a single sense-plus-antisense denominator per dataset.
Genes were taken from the~\syna~annotation (\texttt{syn3a\_genome.gff3}), parsing both the \texttt{gene} and \texttt{pseudogene} feature types so that the three annotated pseudogenes were retained, for a total of 496 loci.
TPM was computed independently from the Illumina and the ONT depth tracks; the single biological replicate of each platform required no replicate averaging, and the two profiles were merged into one table (\texttt{syn3A\_TPM\_Illumina\_ONT.tsv}).
The Illumina profile served as the quantitative transcript-abundance track, with the long-read ONT profile retained as an orthogonal check, and it reproduced the per-gene TPM reported for the same wild-type~\syna~library by~\cite{sandberg_adaptive_2023} (Pearson $r = 0.998$ on $\log_{10}$ TPM), validating the depth-based pipeline.
